\section{Conclusion}
The book supplies valuable genealogical data, documents persistent family
resemblance, and raises consequential questions about its causes. Its
attention to measurement is also well placed. These contributions do not
depend on accepting a single explanation of the correlations.

The genetic interpretation requires more than close fit. The mating
mechanism must generate the equations being estimated, the measurement
model must connect those equations to the recorded outcomes, and the
identifying exclusions must be confronted with evidence. Environmental
comparisons can help, but only against alternatives that predict different
results. Modern genomic designs are especially valuable because random
transmission supplies information unavailable in average kinship
correlations. They can support the additive approximation while challenging
other exclusions.

Even valid genetic estimates would not divide social and genetic causation
into mutually exclusive mechanisms or establish the effects and value of
interventions. Clark recognizes parts of this argument, but his stronger
historical and policy conclusions do not consistently respect it. The
review's disagreement is therefore with the steps connecting evidence to
causal explanation and policy, not with the possibility that genes matter.
